\begin{ExerciseBlock}
\ExerciseHeading{03}{De-Broglie-Wellenlängen}{quantum_optics_de_broglie_wavelengths}
Vergleichen Sie die De-Broglie-Wellenlängen eines Elektrons und eines makroskopischen Körpers.
\begin{enumerate}
\item Ein zunächst ruhendes Elektron wird durch $U=\qty{1.0}{\kilo\volt}$ beschleunigt. Bestimmen Sie seine Geschwindigkeit $v_\mathrm e$ und seine De-Broglie-Wellenlänge $\lambda_\mathrm e$. Verwenden Sie die nichtrelativistische Näherung.
\item Ein Fußball der Masse $m_\mathrm F=\qty{0.45}{\kilogram}$ bewegt sich mit $v_\mathrm F=\qty{30}{\metre\per\second}$. Berechnen Sie seine De-Broglie-Wellenlänge.
\item Vergleichen Sie die Ergebnisse. Bei welchen räumlichen Strukturen sind Welleneffekte beim Elektron relevant? Warum spielen sie beim Fußball praktisch keine Rolle?
\end{enumerate}
\end{ExerciseBlock}
\begin{Solution}
\Needspace{4\baselineskip}
\par\medskip\noindent\textbf{Musterlösung}\par\nopagebreak[4]
\begin{enumerate}
\item Aus $eU=m_\mathrm e v_\mathrm e^2/2$ und $p_\mathrm e=m_\mathrm e v_\mathrm e$ folgt
\[
 v_\mathrm e=\sqrt{\frac{2eU}{m_\mathrm e}}\approx\qty{1.88e7}{\metre\per\second},
 \qquad
 \lambda_\mathrm e=\frac{h}{\sqrt{2m_\mathrm e eU}}\approx\qty{3.88e-11}{\metre}.
\]
\item Für den Fußball ist $\lambda_\mathrm F=h/(m_\mathrm Fv_\mathrm F)\approx\qty{4.9e-35}{\metre}$.
\item Die Wellenlänge des Elektrons liegt in der Größenordnung atomarer Abstände. Dadurch können etwa Kristallgitter Elektronen beugen. Die Wellenlänge des Fußballs ist für das Experiment verschwindend klein.
\end{enumerate}
\end{Solution}
